\documentclass[a4paper,12pt]{book} % nie: report!


\usepackage[T1,plmath]{polski} % lepiej to zamiast babel!
\usepackage[utf8]{inputenc} % w razie kłopotów spróbować: \usepackage[utf8x]{inputenc}
\usepackage{fancyhdr} % nagłówki i stopki
\usepackage{indentfirst} % WAŻNE, MA BYĆ!
\usepackage[pdftex]{graphicx} % to do wstawiania rysunków
\usepackage{amsfonts} % pakiety od AMS, ułatwiają składanie pewnych techniczno-matematcyznych rzeczy
\usepackage{amsmath} % to do dodatkowych symboli, przydatne
\usepackage{amssymb} % to też do dodatkowych symboli, też przydatne
\usepackage{amsthm}
\usepackage[pdftex,
            left=1.1in,right=1.1in,
            top=1.1in,bottom=1.1in]{geometry} % marginesy
\usepackage{float}
\usepackage[font=small,labelfont=bf]{caption}

\usepackage[colorlinks=true]{hyperref} % odnośniki interaktywne w PDFie
\hypersetup{allcolors=black}

\usepackage{listings}
\lstset{
    basicstyle=\footnotesize\tt,
    numbers=left,
    numberstyle=\tiny,
    frame=tb,
    tabsize=4,
    columns=fixed,
    showstringspaces=false,
    showtabs=false,
    keepspaces,
    commentstyle=\color{red},
    keywordstyle=\color{blue}
}
\newfloat{lstfloat}{htbp}{lolst}[chapter]
\floatname{lstfloat}{Listing}
\def\lstfloatautorefname{Listing}

% jesli potrzeb, można oczywiście wstawić inne pakiety i swoje definicje...



% definicje nagłówków i stopek
\pagestyle{fancy}
\renewcommand{\chaptermark}[1]{\markboth{#1}{}}
\renewcommand{\sectionmark}[1]{\markright{\thesection\ #1}}
\fancyhf{}
\fancyhead[LE,RO]{\footnotesize\bfseries\thepage}
\fancyhead[LO]{\footnotesize\rightmark}
\fancyhead[RE]{\footnotesize\leftmark}
\renewcommand{\headrulewidth}{0.5pt}
\renewcommand{\footrulewidth}{0pt}
\addtolength{\headheight}{1.5pt}
\fancypagestyle{plain}{\fancyhead{}\cfoot{\footnotesize\bfseries\thepage}\renewcommand{\headrulewidth}{0pt}}


% interlinia
\linespread{1.25}



\begin{document}
\begin{titlepage}
    ~

    \begin{tabular}{c@{\hspace{21mm}}|@{\hspace{5mm}}l}
        \vspace{-20mm} &                                                                                   \\
        \multicolumn{2}{l}{\hspace{-12.5mm} \includegraphics[width=8cm]{LogoUMCS.jpg}}                     \\
        \multicolumn{2}{@{\hspace{20mm}}l}{\vspace{-4mm}}                                                  \\
        \multicolumn{2}{@{\hspace{28mm}}l}{\Large \sf UNIWERSYTET MARII
        CURIE-SKŁODOWSKIEJ}                                                                                \\
        \multicolumn{2}{@{\hspace{28mm}}l}{\vspace{-4mm}}                                                  \\
        \multicolumn{2}{@{\hspace{28mm}}l}{\Large \sf W LUBLINIE}                                          \\
        \multicolumn{2}{@{\hspace{28mm}}l}{\vspace{-4mm}}                                                  \\
        \multicolumn{2}{@{\hspace{28mm}}l}{\Large \sf Wydział Matematyki, Fizyki i
        Informatyki}                                                                                       \\
        \multicolumn{2}{@{\hspace{28mm}}l}{\vspace{21mm}}                                                  \\
                       & {\sf Kierunek: \textbf{informatyka} }                                             \\
        %& {\sf Specjalność: \textbf{informatyczna}} \\ % wpisujemy tylko jeśli jest!!!
                       &                                                                                   \\\\\\
                       & {\sf \large \bfseries Jan Bylina}                                                 \\
                       & {\sf nr albumu: 303827}                                                           \\
                       &                                                                                   \\\\\\
                       & \Large \sf \bfseries Projekt oraz implementacja systemu                           \\
                       & \Large \sf \bfseries gromadzenia rozproszonych danych                             \\
                       & \Large \sf \bfseries z wykorzystaniem technologii LoRa                            \\\\[-10pt]
                       & {\large \sf Design and implementation of the distributed data collection system } \\
                       & {\large \sf using LoRa technology}                                                \\
                       &                                                                                   \\
                       &                                                                                   \\
                       &                                                                                   \\
                       & {\sf Praca licencjacka}                                                           \\
                       & \vspace{-7mm}                                                                     \\
                       & {\sf napisana w Katedrze Oprogramowania Systemów Informatycznych}                 \\
                       & {\sf Instytutu Informatyki UMCS}                                                  \\
                       & \vspace{-7mm}                                                                     \\
                       & {\sf pod kierunkiem \bfseries dr hab. Przemysława Stpiczyńskiego}                 \\
        \multicolumn{2}{@{\hspace{28mm}}l}{\vspace{15mm}}                                                  \\
        \multicolumn{2}{@{\hspace{28mm}}l}{\textbf{\textsf{Lublin 2023}}}
    \end{tabular}
\end{titlepage}




\sloppy



\thispagestyle{empty}


\newpage{}

\thispagestyle{empty}

\newpage{}



\tableofcontents{}

\chapter*{Wstęp} % z gwiazdką, więc bez numerka...
\addcontentsline{toc}{chapter}{Wstęp} % ...ale w spisie treści ma być

Tu treść wstępu
WSTĘP
WSTEP
------

\chapter{Rozdział — tutorial}

\section{Sekcja A}

W tabeli~\ref{tab:przyk} widzimy przykład tabeli z~nagłówkiem i~odnośnikiem. Tabele tworzymy z~nagłówkiem na górze oraz opcją \texttt{[t]}.
Natomiast na rysunku~\ref{rys:przyk} --- widzimy przykład rysunku z nagłówkiem i~odnośnikiem.
Rysunki tworzymy z nagłówkiem pod spodem oraz opcją \texttt{[b]}.
Rysunki powinny być w formacie PDF; jeśli to niemożliwe, to PNG (w wysokiej rozdzielczości); a~ostatecznie JPG (jak tu). Jeśli chcemy sterować rozmiarem, to zwykle najwygodniej użyć \texttt{width=...}
Ponadto możemy odwoływać się do bibliografii~\cite{bib:a, bib:b}.

Jeśli chodzi o wzory, możemy złożyć je na kilka sposobów, w zależności od potrzeb --- w tekście: $e=\lim_{n\to\infty}\left(1+\frac{1}{n}\right)^n$, wyniesiony do osbnej linii
(warto zwrócić uwagę, że ten i~kolejny są złożone nieco inaczej niż pierwszy):
\[e=\lim_{n\to\infty}\left(1+\frac{1}{n}\right)^n,\] a także wyniesiony z numerem:
\begin{equation}
    e=\lim_{n\to\infty}\left(1+\frac{1}{n}\right)^n.
    \label{wzor:e}
\end{equation}
Do tego oostatniego możemy się odwołać:~\eqref{wzor:e}.

No i oczywiście listingi --- listing~\ref{lst:przyk} pokazuje, jak zrobić to w~miarę poprawnie\ldots{}

\begin{table}[t]
    \begin{center}
        \caption{Przykładowa tabela}\label{tab:przyk}
        \begin{tabular}{l|c|r}
            slkdjfslj  & sdkskd               & s;lkdsdk          \\
            \hline
            slkjd      & skljdsldj            & skljdsjdsldj      \\
            sljkdslkjd & woieupowiepoweiwiewp & weoiw eppowie wpo \\
        \end{tabular}
    \end{center}
\end{table}

\begin{figure}[b]
    \begin{center}
        \includegraphics[width=3cm]{LogoUMCS}
    \end{center}
    \caption{Przykładowy rysunek}\label{rys:przyk}
\end{figure}


\begin{lstfloat}[b]
    \lstset{language=C++}
    \begin{lstlisting}[frame=single]
tab[0:n] = dem[nRows][nCols]; //?
#pragma acc data copy(tab [0:n], slope [0:n])
\end{lstlisting}
    \caption{Jakieś dwie linijki w~C++ (z~OpenACC)}\label{lst:przyk}
\end{lstfloat}
\
sekcja 232323
\

\section{Sekcja B}

sekcja BB

\chapter{Wykorzystane narzędzia, technologie i protokoły}

\section{Urządzenia}
\subsection{STM32}
---
\subsection{ESP32}
---
\subsection{Rasberry Pi Pico}
---

\section{Języki programowania i technologie}
\subsection{\texttt{C++} for Arduio}
---
\subsection{\texttt{C} for STM32}
---
\subsection{\texttt{MicroPython} for Rasberry Pi Pico}
---
\subsection{\texttt{Python} for MQTT}

\section{Protokoły komunikacyjne}
---
\subsection{\texttt{MQTT}}
---
\subsection{\texttt{LoRa}}
---
\subsection{\texttt{HTTP}}
---
\section{Bazy danych i pozostałe technologie}
\subsection{InfluxDB 2}
---
\subsection{Docker}
---
\subsection{PlatformIO}
---
\chapter{Istniejące rozwiązania}
---
\section{LoRaWAN}
---
\subsection{The Things Network ?}
---
\subsection{ChirpStack ?}
---
\subsection{Loriot ?}
---
\section{Artykuły}
---
\section{Wpisy w sieci i blogach}

\chapter{Założenie i Implementacja}
\chapter{Wdrożenie i testy}
\chapter{Wnioski i perspektywy rozwoju}



\listof{lstfloat}{Spis listingów} % jeśli są listingi
\addcontentsline{toc}{chapter}{Spis listingów}

\listoftables{} % jeśli są tabele
\addcontentsline{toc}{chapter}{Spis tabel}

\listoffigures{} % jeśli są rysunki
\addcontentsline{toc}{chapter}{Spis rysunków}

\begin{thebibliography}{99}
    \addcontentsline{toc}{chapter}{Bibliografia} % też ręczne dodanie do spisu treści, jak Wstęp
    \bibitem{bib:a} aaaaaaaa
    \bibitem{bib:b} bbbbbbbb
\end{thebibliography}


\end{document}
